% =====================================================================
% P5 merge snippet — M06 / FM03: contributions table
% Source: manuscript_v2, Sec. 1 (tab:contributions, C1--C7), written 2026-11-01.
% Adapted: every row now points at a section or appendix of THIS paper, and every evidence
% cell cites only artefacts that exist in this paper (app:registry's 14 benchmarks,
% app:repro's scripts, app:params' parameter register, app:disp's derivation,
% sec:synthesis' comparison tables). Nothing is claimed that this paper cannot show.
% NEW LABELS: tab:contributions
% Suggested placement: \cref{sec:intro_ref}, immediately before the outline of the paper.
% =====================================================================

The contributions a reader can cite this review for are collected in \cref{tab:contributions}. Each row
names the place where the contribution is delivered, so that the claim can be checked rather than
believed, and the evidence column names only what this paper actually contains.

\begin{table}[H]\centering\small
\caption{What this review contributes, where it is delivered, and what supports it.}
\label{tab:contributions}
\begin{tabularx}{\textwidth}{@{}p{0.5cm}Y p{2.5cm}Y@{}}
\toprule
\textbf{\#} & \textbf{Contribution} & \textbf{Delivered in} & \textbf{Evidence in this paper} \\
\midrule
C1 & One heat-law template and its stepwise reductions, so that LS, GL, GN-II/III, DPL, TPL, MGT, MDD
  and fractional models are compared as members of one family rather than as competing names. &
  \cref{sec:foundations}, \cref{sec:hyperbolic}, \cref{sec:lags}, \cref{sec:memory} &
  \eqref{eq:template} and the reduction table; each reduction checked symbolically \\
C2 & A positive-definite availability functional with its flux, the dissipation identity it satisfies,
  and the explicit statement that it is not the total-energy flux --- together with the exact Eulerian
  total-energy identity and its first-order reduction. & \cref{sec:foundations}, \cref{sec:interfaces} &
  \eqref{eq:energy}, \eqref{eq:P}, \eqref{eq:exactflux}; energy ratios summing to one to machine precision \\
C3 & A family-comparative synthesis with a decision guide: what each theory predicts, what it costs, and
  under which conditions two families may legitimately be compared. & \cref{sec:synthesis} &
  comparison tables in \cref{sec:synthesis}; one worked example per family in \cref{sec:examples} \\
C4 & An original five-axis taxonomy that places any candidate model, hybrids included, on one set of
  coordinates, with an explicit dissimilarity that says when two studies are comparable at all. &
  \cref{sec:taxonomy} & axis definitions, the dissimilarity \eqref{eq:taxdissim}, and the assignments in
  \cref{tab:taxonomy-axes} \\
C5 & A per-equation verification register: every registered equation carries a status from a finite
  vocabulary, so that derived, primary-source-checked, unverified and open are distinguishable. &
  \cref{sec:eqregister} & \cref{tab:eqregister}; four entries at PRIMARY-SOURCE-CHECKED with printed page
  and equation \\
C6 & Fourteen executable benchmarks with reference solutions and stated tolerances, each a code
  verification rather than a physical validation. & \cref{app:registry}, \cref{app:repro} &
  \cref{app:registry} (14 rows), the scripts of \cref{app:repro}, and the parameter register
  \cref{app:params} \\
C7 & A stated evidence protocol and a set of reporting standards, so that a reader can discount each
  statement by the rule it fails and a future author can make a result comparable. &
  \cref{sec:evidence_protocol}, \cref{sec:reporting} & the six rules, and the reporting checklist \\
\bottomrule
\end{tabularx}
\end{table}

One caution belongs with C5, because it is the row a reader is most likely to over-read: a register entry
records the status of a check, not the importance or the validity of an equation for a given material.
Mechanical correctness does not make a constitutive assumption true.
